% Pursuing Stacks (cotes 134-2 to 134-8): archive cote 134-3, site page 14 (PDF page 15), typescript page 14.
% typescript. Draft compared: ps2.tex:616-682.
% AI-transcribed by Claude Opus 5.5 (claude-opus-5-5, adaptive thinking, effort high), Message Batches API, batch ps1 (msgbatch_01Rv5hCaNTtUt4JHx1yqoHPL), 2026-11-12; not yet checked by eye.
\scanpage{134-3/14}
\hand{5}.3. \mspage{\hand{14}}

\hand{20.} I had barely st opped writing last tuesday,  when it became clear \xout{th}
that the ``quick glance'' had been too quick indeed. \xout{It would} As a matter
of fact,  the ``topology'' I was contemplating on $\underline{B}$ in terms of covering
families does not satisfy the conditions for a ``site'' -  that something
was fishy first occured to me through the heuristic consequence,  that
the functors  ``n-th component''
\[ F_{*} \hand{$\longmapsto$} F_{n} \]
from \xout{\illegible{1}} stacks to sets are fiber functors,  or equivalently, that
direct limits in the category of Gr-stacks can be computed componentwise
-  now this is definitely false,  for much the same reasons why it is
false \unsure{already|alfeady} in the category (Cat) of categories.  However,  it occured
to me that \inserted{in} the latter category (Cat),  although not a topos (for instan-%
ce because, according to Giraud's paper on descent theory of 1965\hand{, (19)}
the implications for epimorphisms
\diagram{
\begin{tikzcd}[row sep=tiny,column sep=large]
 & \text{effective} \arrow[dr, Rightarrow, "\text{(hand)}"'] & \\
\text{effective universal} \arrow[ur, Rightarrow] \arrow[dr, Rightarrow] & & \text{just epi} \\
 & \text{universal} \arrow[ur, Rightarrow] &
\end{tikzcd}
}
are strict),  there \uline{is} a very natural topology,  \xout{namely} turning it into
a site,  namely the one where a family of morphisms namely functors
\[ A_{i} \hand{$\longrightarrow$} A \]
is covering if and only if the corresponding family in (Ss sets)
\[ \mathrm{Nerve}(A_{i}) \hand{$\longrightarrow$} \mathrm{Nerve}(A) \]
is covering,  i.e. iff every \xout{family} sequence of composable arrows in A
\[ a_{0} \hand{$\to$} a_{1} \hand{$\to$} a_{2} \hand{$\to$} a_{3} \hand{$\to$} \quad \dots \hand{$\to$} a_{n} \]
can be lifted  to one among the $A_{i}$'s.  As a matter of fact,  this condi-%
\margin{5.3. This is definitely false, see \S\unsure{24} below}
tion (where it suffices to take n=2,  visibly)  is equivalent (according
to Giraud) to the condition that the family be \hand{``}universally effectively
epimorphic\hand{''},  i.e. covering with respect to the``canonical topology''of
(Cat)\hand{(15)}.This suggests a third way for associating to a \inserted{[small],} category A a
topology-like structure,  namely the topos of all sheaves over
$(\mathrm{Cat})_{/A}$  ,  the site of all categories over A, \inserted[typed]{endowed} with the topology induced
by the canonical topology of (Cat)  (which is indeed the canonical topo\hand{-}%
logy of  $(\mathrm{Cat})_{/A}$ \hand{)}.  It should be an easy exercice in terms of nerves to
check that the homotopy type (a priori,  a prohomotopy type)  associated
to this site is just \hand{``}the\hand{''}homotopy type of A,  defined either as the
homotopy type of Nerve(A),  or as the (pro)homotopy type of the topos
$\hand{\hat{A}}$ of all presheaves over A.  This of course parallels the similar fami-%
liar fact in the category (Ss sets) of all semisimplicial sets,  namely
that for such a ss set\deleted{\unsure{s}} K,  the homotopy type of K can be viewed also
as the homotopy type of the induce\hand{d} topos  $(\mathrm{Ss\ sets})_{/K}$ of all ss sets
over K.  The difference in the two cases is that  in the second case
